\section{Strongly regular graphs}
Let $\Gamma$ be noncomplete and nonempty, strongly regular with parameters
$(N,k,\lambda,\mu)$. The three kinds of adjacency entries give
\[
 A^2=(k-\mu)I+(\lambda-\mu)A+\mu J.
\]
Thus \Uref{ULCD.SRG} recovers both parts of source Proposition~4:
\[
\begin{split}
 C\text{ is RLCD}&\iff k=0,\quad\lambda=-1,\quad\mu=0\pmod p,\\
 C\perp\bar C&\iff k-\mu=0,\quad k-1-\lambda=0\pmod p.
\end{split}
\]
In the second line CLCD holds exactly when $k$ and $N-1-k$ do not both
vanish. This completes the possible-exception clause in the source.
In the first line the nonconstant eigenvalues, when integral, have
residues zero and minus one: their product is $\mu-k=0$ and their sum
is $\lambda-\mu=-1$. Hence the dimension is the multiplicity of the
minus-one residue. The radical in the usual formula is
\[
 \Delta=(\lambda-\mu)^2+4(k-\mu),
\]
with a plus sign. The minus sign printed in Proposition~4(1) is a
transcription error, inconsistent with the preceding source spectrum.
The classical multiplicities used here are
\[
 m_\pm=\frac12\left(N-1\pm
 \frac{(N-1)(\mu-\lambda)-2k}{\sqrt\Delta}\right),
 \qquad \rho_\pm=\frac{\lambda-\mu\pm\sqrt\Delta}{2}.
\]
They follow from $m_++m_-=N-1$ and $k+m_+\rho_++m_-\rho_-=0$.
The source's cited conference-parameter alternative has
$\lambda=\mu-1$, $k=2\mu$, $N=4\mu+1$ and $N$ a sum of two
integer squares; the other spectral alternative has integral eigenvalues.
That classical feasibility statement is retained from the source's
Cameron--van Lint reference, not reproved here or needed to test the
displayed graph families. On each nonconstant eigenspace the complement
eigenvalue is $-1-\rho$, with the same multiplicity.
In particular $T(m)$ has nonconstant eigenvalues $m-4,-2$, of
multiplicities $m-1,m(m-3)/2$, and its complement has $3-m,1$ on
those respective spaces. This corrects the mismatched complement
entries in the source's triangular spectral display.

\begin{proposition}\label{prop:srg-applications}
The four instances in source Corollary~2 have the following conclusions.
\begin{enumerate}
\item For odd $m\geq5$, the binary triangular graph $T(m)=J(m,2)$
is RLCD, of dimension $m-1$. It is CLCD exactly when $m\equiv1\pmod4$.
Its complement $KG(m,2)$ is RLCD exactly for $p=2$ and $m\equiv3\pmod4$.
\item For odd $m\geq3$, the binary lattice graph $L_2(m)$ and its
complement are both RLCD, of dimensions $2m-2$ and $(m-1)^2$.
\item Let $q\equiv1\pmod4$ be any odd prime power. The Paley graph
$P(q)$ is RLCD over $\mathbb F_p$ exactly when $p=2$ and $q\equiv1\pmod8$,
or when $p$ is odd and $q\equiv1\pmod p$. In each case the dimension is
$(q-1)/2$, including when $q$ is not a square.
\item Every strongly regular graph with parameters $(81,20,1,6)$ is
binary RLCD of dimension $20$. The specified generalized Paley graph
below realizes the source example.
\end{enumerate}
\end{proposition}
\begin{proof}
For $T(m)$, the parameters are
$(\binom m2,2(m-2),m-2,4)$. The reflexive test forces $p=2$ and odd $m$.
The simplex theorem gives the dimension. In that characteristic its
degree vanishes and the exact defect theorem gives CLCD precisely when
$\binom m2$ is even, namely $m\equiv1\pmod4$. Substituting the
complement parameters gives the final condition; this also follows
from the complete subset classification in the common paper.

For $L_2(m)$, the parameters are $(m^2,2(m-1),m-2,2)$, and the test gives
$p=2$ and odd $m$. Its nonconstant eigenvalues are $m-2$ and $-2$, of
multiplicities $2(m-1)$ and $(m-1)^2$. The LCD rank lemma gives the first
dimension. Equation~\eqref{eq:two-reflexive} proves that the complement
is also RLCD. Both degrees vanish modulo two, so the exact defect gives
its dimension $m^2-1-(2m-2)=(m-1)^2$.

Paley parameters are $(q,(q-1)/2,(q-5)/4,(q-1)/4)$, so the displayed
conditions are exactly the reflexive test. The graph is self-complementary:
multiplication of $\mathbb F_q$ by any nonsquare interchanges square and
nonsquare differences and hence the two adjacency matrices. Their ranks
are equal over every field. In each stated characteristic both degrees
vanish, so the constant-line defect gives
$2\rank A=q-1$. This proves the dimension without integral eigenvalues.

For $(81,20,1,6)$, the binary reflexive test holds. The two nonconstant
eigenvalues are $2$ and $-7$, of multiplicities $60$ and $20$. Only the
latter is odd, and the rank lemma proves the final assertion.
\end{proof}

The displayed triangular exceptional case in the source is genuine:
$m\equiv3\pmod4$ gives both degrees zero and a one-dimensional missing
constant line. The needed congruence is $\binom m2\equiv1\pmod2$;
the stronger modulus-four assertion printed there does not hold for
every such $m$. The binary Petersen graph is not RLCD, but, as explained
by the subset classification, it is CLCD.

For the symplectic graph on projective points of a $2m$-dimensional
symplectic space, $m\geq2$, put
\[
 N=\frac{q^{2m}-1}{q-1},\quad
 k=\frac{q^{2m-1}-1}{q-1}-1,\quad
 \mu=\frac{q^{2m-2}-1}{q-1},\quad\lambda=\mu-2.
\]
No field characteristic can make this graph RLCD: the requirements
$\mu=0$ and $\lambda=-1$ would give $-2=-1$. Nor can it be CLCD, since
$k-1-\lambda=(k-\mu)+1$ makes the two complement tests incompatible.
For its complement the adjacent and nonadjacent intersection numbers
are equal, so its reflexive tests are also incompatible. These direct
parameter arguments establish all exclusions in the source.
For completeness the nonconstant eigenvalues are $q^{m-1}-1$ and
$-q^{m-1}-1$, with multiplicities
\[
 \tfrac12(N-1+q^m),\qquad \tfrac12(N-1-q^m).
\]
The source prints $\mu-1$ instead of $N-1$ in both multiplicities;
the correction is forced by their sum $N-1$ and the trace equation.
No erroneous spectral multiplicity is needed for the exclusion proof.
